% =============================================================================
% Resumen teorico - Practico 03 (Autoreferencia y Teorema de la Recursion)
% Computabilidad y Complejidad (Teoria de la Computacion II)
% FCEFQyN - Universidad Nacional de Rio Cuarto
% Primer Cuatrimestre de 2026
% -----------------------------------------------------------------------------
% Fuentes principales:
%   - chapter-06.pdf  (Sipser, Cap. 6: Advanced Topics in Computability Theory)
%   - autoref.pdf     (laminas de la catedra: AutoReferencia)
%
% Compilacion sugerida: tectonic resumen-teorico-practico-03.tex
% =============================================================================

\documentclass[11pt,a4paper]{article}

% --- Codificacion y idioma --------------------------------------------------
\usepackage[utf8]{inputenc}
\usepackage[T1]{fontenc}
\usepackage[spanish,es-tabla]{babel}
\usepackage{lmodern}
\usepackage{microtype}

% --- Matematica -------------------------------------------------------------
\usepackage{amsmath,amssymb,amsthm,mathtools}
\usepackage{stmaryrd}

% --- Layout y tipografia ----------------------------------------------------
\usepackage[margin=2.7cm]{geometry}
\usepackage{enumitem}
\setlist[itemize]{topsep=2pt,itemsep=2pt,parsep=0pt}
\setlist[enumerate]{topsep=2pt,itemsep=2pt,parsep=0pt}
\usepackage{parskip}
\usepackage{xcolor}
\usepackage{booktabs}
\usepackage{array}

% --- Codigo -----------------------------------------------------------------
\usepackage{fancyvrb}
\fvset{
  fontsize=\footnotesize,
  frame=single,
  framerule=0.3pt,
  xleftmargin=6pt,
  xrightmargin=6pt
}
\DefineVerbatimEnvironment{pycode}{Verbatim}{}

% --- Hyperref ---------------------------------------------------------------
\usepackage[hidelinks]{hyperref}
\hypersetup{
  pdftitle={Resumen teorico - Practico 03},
  pdfauthor={Computabilidad y Complejidad - UNRC 2026},
  pdfsubject={Autoreferencia y Teorema de la Recursion - Resumen teorico},
  pdfkeywords={recursion, autoreferencia, SELF, MIN, oraculo, Turing reducibility}
}

% --- Ambientes de teoremas --------------------------------------------------
\theoremstyle{plain}
\newtheorem{teorema}{Teorema}[section]
\newtheorem{proposicion}[teorema]{Proposicion}
\newtheorem{lema}[teorema]{Lema}
\newtheorem{corolario}[teorema]{Corolario}

\theoremstyle{definition}
\newtheorem{definicion}[teorema]{Definicion}
\newtheorem{ejemplo}[teorema]{Ejemplo}

\theoremstyle{remark}
\newtheorem*{observacion}{Observacion}
\newtheorem*{idea}{Idea}
\newtheorem*{intuicion}{Intuicion}

% --- Comandos de notacion ---------------------------------------------------
\newcommand{\MT}{\mathrm{MT}}
\newcommand{\TM}{\mathrm{TM}}
\newcommand{\enc}[1]{\langle #1\rangle}
\newcommand{\encMw}{\langle M, w\rangle}
\newcommand{\lang}{\mathcal{L}}
\newcommand{\Sstar}{\Sigma^{*}}
\newcommand{\redm}{\le_{m}}
\newcommand{\redT}{\le_{T}}
\newcommand{\compl}[1]{\overline{#1}}
\newcommand{\qaccept}{q_{\text{accept}}}
\newcommand{\qreject}{q_{\text{reject}}}
\newcommand{\ATM}{A_{\TM}}
\newcommand{\HALT}{\mathit{HALT}_{\TM}}
\newcommand{\ETM}{E_{\TM}}
\newcommand{\EQTM}{\mathit{EQ}_{\TM}}
\newcommand{\MIN}{\mathit{MIN}_{\TM}}
\newcommand{\SELF}{\mathrm{SELF}}
\newcommand{\Pw}[1]{P_{#1}}
\newcommand{\Powerset}{\mathcal{P}}

% --- Cabecera ---------------------------------------------------------------
\title{\textbf{Resumen te\'orico -- Pr\'actico 03}\\[0.2em]
       \textbf{Autoreferencia, Teorema de la Recursi\'on y Or\'aculos}\\[0.3em]
       \normalsize Computabilidad y Complejidad (Teor\'ia de la Computaci\'on II)\\
       \normalsize FCEFQyN -- Universidad Nacional de R\'io Cuarto\\
       \normalsize Primer Cuatrimestre de 2026}
\author{}
\date{}

\begin{document}
\maketitle
\thispagestyle{empty}

\hrule
\vspace{0.4em}
\begin{center}
\small
Resumen te\'orico-pr\'actico que acompa\~na al Pr\'actico 03. Fuentes
principales: Cap\'itulo 6 del libro de M.~Sipser, \emph{Introduction to the
Theory of Computation} (\texttt{chapter-06.pdf}), y l\'aminas de la c\'atedra
(\texttt{autoref.pdf}).
\end{center}
\vspace{0.4em}
\hrule

\tableofcontents
\bigskip\hrule\bigskip

% =============================================================================
\section{Motivaci\'on: la paradoja de la autoreproducci\'on}

Sipser abre el Cap\'itulo 6 con una vieja paradoja: si toda m\'aquina $A$
que produce otra m\'aquina $B$ es necesariamente \emph{m\'as compleja} que
$B$, entonces ninguna m\'aquina podr\'ia construirse a s\'i misma. Como
nosotros mismos somos m\'aquinas (en alg\'un sentido biol\'ogico) y nos
reproducimos, la paradoja debe estar mal planteada.

El \textbf{Teorema de la Recursi\'on} resuelve la paradoja mostrando que las
m\'aquinas \emph{s\'i} pueden autoreproducirse, e incluso pueden obtener su
propia descripci\'on y operar con ella. Esa capacidad es la base de muchas
herramientas: virus inform\'aticos, quines, introspecci\'on en lenguajes
modernos, y demostraciones limpias de indecibilidad.

\bigskip\hrule\bigskip

% =============================================================================
\section{La funci\'on $q$ y la m\'aquina $\SELF$}

\subsection{El lema constructivo}

\begin{lema}[Sipser 6.1]
\label{lem:q}
Existe una funci\'on computable
\[
  q:\Sstar \longrightarrow \Sstar
\]
tal que, para toda $w\in\Sstar$, $q(w)=\enc{\Pw{w}}$, donde $\Pw{w}$ es una
M\'aquina de Turing que, en cualquier entrada, borra la cinta, escribe $w$ y
se detiene.
\end{lema}

\begin{intuicion}
Dada una cadena $w$, construir ``a mano'' una MT que la imprima es trivial:
cableamos $w$ en una tabla de transiciones que la escribe letra por letra.
Lo no obvio del lema es que esa construcci\'on es \emph{computable} en $w$:
hay una MT $Q$ que, recibiendo $w$, produce la descripci\'on $\enc{\Pw{w}}$.
\end{intuicion}

\subsection{La m\'aquina $\SELF$}

\begin{teorema}[Sipser, p\'aginas 246--248]
Existe una M\'aquina de Turing $\SELF$ que, en cualquier entrada, imprime su
propia descripci\'on $\enc{\SELF}$ en la cinta y se detiene.
\end{teorema}

\begin{proof}[Construcci\'on]
$\SELF$ se descompone en dos partes $A$ y $B$ ejecutadas en secuencia:

\begin{itemize}
  \item \textbf{Parte $A$:} es la m\'aquina $\Pw{\enc{B}}$ del
        Lema~\ref{lem:q}. Al ejecutarse, deja $\enc{B}$ en la cinta.
  \item \textbf{Parte $B$:} con $\enc{B}$ en la cinta:
        \begin{enumerate}
          \item computa $q(\enc{B}) = \enc{A}$;
          \item combina $\enc{A}$ y $\enc{B}$ en $\enc{AB}=\enc{\SELF}$;
          \item imprime $\enc{\SELF}$ y se detiene.
        \end{enumerate}
\end{itemize}

Observar que no hay circularidad: $\enc{A}$ no aparece como dato cableado en
$B$; se \emph{reconstruye en tiempo de ejecuci\'on} aplicando $q$ a la
cadena $\enc{B}$ que dej\'o $A$.
\end{proof}

\subsection{Analog\'ia textual}

Sipser ilustra la t\'ecnica con una oraci\'on autoreproductiva:

\begin{center}\emph{
  Copiar dos veces el siguiente p\'arrafo, una con doble comillas:\\
  ``Copiar dos veces el siguiente p\'arrafo, una con doble comillas:''}
\end{center}

La parte $B$ es la oraci\'on sin comillas; la parte $A$ es la misma oraci\'on
entre comillas. Al ``ejecutar'' las instrucciones se obtiene el texto
original. La l\'inea con comillas es exactamente $\Pw{\enc{B}}$ (un dato
literal); la l\'inea sin comillas es el ``programa'' que la consume.

\bigskip\hrule\bigskip

% =============================================================================
\section{El Teorema de la Recursi\'on}

\begin{teorema}[Sipser 6.3, recursi\'on]
\label{teo:recursion}
Sea $T$ una M\'aquina de Turing que computa una funci\'on
$t: \Sstar\times\Sstar \to \Sstar$. Entonces existe una MT $R$ que computa
una funci\'on $r:\Sstar\to\Sstar$ tal que, para toda $w\in\Sstar$,
\[
  r(w) \;=\; t\bigl(\enc{R},\, w\bigr).
\]
\end{teorema}

\begin{intuicion}
$T$ es ``lo que se quiere hacer'' suponiendo que ya se tiene $\enc{R}$ a
mano; el teorema garantiza que existe una MT $R$ que se comporta como $T$
con su propia descripci\'on autom\'aticamente inyectada en la primera
componente. En otras palabras: una MT puede \emph{obtener su propia
descripci\'on} y computar con ella.
\end{intuicion}

\begin{proof}[Esquema]
Se construye $R$ en tres partes $A,B,T$:
\begin{itemize}
  \item $A=\Pw{\enc{BT}}$: deja $\enc{BT}$ en la cinta sin tocar la entrada
        $w$.
  \item $B$: aplica $q$ a $\enc{BT}$, obtiene $\enc{A}$, combina con $BT$
        para obtener $\enc{ABT}=\enc{R}$, escribe $\enc{R,w}$ y pasa el
        control a $T$.
  \item $T$ termina la computaci\'on como en el enunciado.
\end{itemize}
$R$ resulta as\'i una MT bien definida con la propiedad pedida.
\end{proof}

\subsection*{Diagrama mental}

\begin{center}
\begin{tabular}{ccc}
$\enc{M},\;w$ & $\rightarrow$ & $T(\enc{M},w)$  \\
$w$           & $\rightarrow$ & $R$ usa su propia descripci\'on para computar $T(\enc{R},w)$ \\
\end{tabular}
\end{center}

\subsection*{C\'omo usar el teorema}

A la hora de dise\~nar una MT, se puede simplemente incluir en su pseudoc\'odigo
la frase ``\emph{obtener la propia descripci\'on $\enc{M}$ por el Teorema
de la Recursi\'on}''. El Teorema garantiza que tal m\'aquina existe.

\bigskip\hrule\bigskip

% =============================================================================
\section{Aplicaciones cl\'asicas del Teorema de la Recursi\'on}

\subsection{Indecibilidad de $\ATM$ (nuevo enfoque)}

Una de las consecuencias m\'as elegantes del Teorema de la Recursi\'on es una
prueba muy corta de la indecibilidad de
\[
  \ATM \;=\; \{\encMw \mid M \text{ acepta } w\}.
\]

\begin{teorema}[Sipser 6.5]
$\ATM$ es indecidible.
\end{teorema}

\begin{proof}
Supongamos, para llegar a una contradicci\'on, que existe una MT $H$ que
decide $\ATM$. Definimos:

\medskip\noindent\textbf{$B = $ ``Con entrada $w$:}
\begin{enumerate}
  \item obtener su propia descripci\'on $\enc{B}$ (Teorema de la Recursi\'on);
  \item simular $H$ con entrada $\enc{B,w}$;
  \item hacer lo contrario de lo que $H$ dice: aceptar si $H$ rechaza, y
        rechazar si $H$ acepta.''
\end{enumerate}
\medskip

Para cualquier $w$, $B$ acepta $w$ si y solo si $H$ rechaza $\enc{B,w}$, es
decir, si y solo si $B$ \emph{no} acepta $w$. Contradicci\'on. Luego $H$ no
puede existir.
\end{proof}

\begin{observacion}
La prueba evita la diagonalizaci\'on de Cantor (Teorema 4.11) y reemplaza el
truco de ``aplicarse al propio c\'odigo'' por la herramienta general del
Teorema de la Recursi\'on. El resultado es el mismo, pero el argumento se
vuelve casi mec\'anico.
\end{observacion}

\subsection{Las MT minimales no son reconocibles}

\begin{definicion}[Sipser 6.6]
Si $M$ es una MT, la \emph{longitud} de $\enc{M}$ es la cantidad de s\'imbolos
de la cadena que la describe. Decimos que $M$ es \emph{minimal} si no existe
una MT equivalente con descripci\'on m\'as corta. Definimos
\[
  \MIN \;=\; \{\enc{M} \mid M \text{ es minimal}\}.
\]
\end{definicion}

\begin{teorema}[Sipser 6.7]
$\MIN$ no es Turing-reconocible.
\end{teorema}

\begin{proof}
Supongamos que un enumerador $E$ enumera $\MIN$. Construimos:

\medskip\noindent\textbf{$C = $ ``Con entrada $w$:}
\begin{enumerate}
  \item obtener su propia descripci\'on $\enc{C}$ (Recursi\'on);
  \item ejecutar $E$ hasta que aparezca una MT $D$ con $|\enc{D}|>|\enc{C}|$;
  \item simular $D$ sobre $w$.''
\end{enumerate}
\medskip

Como $\MIN$ es infinito, eventualmente $E$ imprime una MT $D$ con
$|\enc{D}|>|\enc{C}|$. Por construcci\'on, $C$ es \emph{equivalente} a $D$
(simula a $D$ paso a paso). Pero $|\enc{C}|<|\enc{D}|$, luego $D$ \emph{no}
es minimal. Contradicci\'on con $\enc{D}\in\MIN$.
\end{proof}

\subsection{Teorema del punto fijo}

\begin{teorema}[Sipser 6.8, punto fijo]
\label{teo:fixedpoint}
Sea $t:\Sstar\to\Sstar$ computable. Existe una MT $F$ tal que $t(\enc{F})$
describe una MT equivalente a $F$.
\end{teorema}

\begin{proof}[Esquema]
Definimos $F$:
\begin{enumerate}
  \item Obtener $\enc{F}$ (Recursi\'on).
  \item Computar $t(\enc{F})$, que es la descripci\'on de cierta MT $G$.
  \item Simular $G$ sobre la entrada.
\end{enumerate}
Entonces $F$ y $G$ tienen comportamientos id\'enticos, y $t(\enc{F})=\enc{G}$
describe una MT equivalente a $F$.
\end{proof}

El teorema dice que toda transformaci\'on computable de descripciones de MT
tiene un \emph{punto fijo sem\'antico}: alguna MT cuyo comportamiento queda
inalterado por la transformaci\'on. Es el an\'alogo del Teorema de Rice en
su forma m\'as primitiva.

\bigskip\hrule\bigskip

% =============================================================================
\section{Introspecci\'on en lenguajes de programaci\'on}

El Teorema de la Recursi\'on garantiza la posibilidad de \emph{obtener la
propia descripci\'on}. En lenguajes de programaci\'on reales esta capacidad
se llama \textbf{introspecci\'on} (o \emph{reflection}).

\subsection{El m\'odulo \texttt{inspect} de Python}

Python expone esta capacidad mediante:
\begin{itemize}
  \item \texttt{inspect.getsource(obj)}: devuelve el c\'odigo fuente de
        \texttt{obj} (funci\'on, clase, etc.). Es la versi\'on real de
        ``obtener $\enc{M}$''.
  \item \texttt{inspect.getmembers(obj)}: lista todos los miembros de un
        objeto: atributos, m\'etodos, dunders.
  \item \texttt{inspect.signature(funci\'on)}: devuelve la firma
        (par\'ametros y anotaciones).
  \item \texttt{inspect.isfunction}, \texttt{inspect.isclass}, etc.:
        predicados para clasificar miembros.
\end{itemize}

Ejemplo m\'inimo:

\begin{pycode}
import inspect
def saludar(nombre: str) -> str:
    return f"Hola, {nombre}"
print(inspect.getsource(saludar))
print(inspect.signature(saludar))
\end{pycode}

\subsection{Quines: el patron $A+B$}

Un \emph{quine} es un programa que imprime su propio codigo fuente. La forma
canonica usa el patron $A+B$ de $\SELF$. En Python:

\begin{pycode}
B = 'B = {!r}\nprint(B.format(B))\n'
print(B.format(B))
\end{pycode}

Aqui \texttt{repr} cumple el rol de la funcion $q$ del Lema~\ref{lem:q}:
toma una cadena y produce el literal Python que la representa.

\subsection{Creaci\'on din\'amica de clases}

Python permite construir clases en tiempo de ejecucion con la built-in
\texttt{type(nombre, bases, atributos)}. Es exactamente el constructor que
invoca internamente la sentencia \texttt{class}:

\begin{pycode}
Punto = type("Punto", (object,), {
    "__init__": lambda self, x, y: setattr(self, "x", x) or setattr(self, "y", y),
    "norma": lambda self: (self.x ** 2 + self.y ** 2) ** 0.5,
})
\end{pycode}

Esto es la version programable de la construccion del Lema~\ref{lem:q}: a
partir de una \emph{descripcion} (nombre, bases, diccionario) se produce un
nuevo objeto ejecutable.

\bigskip\hrule\bigskip

% =============================================================================
\section{M\'aquinas con or\'aculo y reducibilidad de Turing}

\subsection{Definiciones}

\begin{definicion}[Sipser 6.18, or\'aculo]
Un \emph{or\'aculo} para el lenguaje $B$ es un dispositivo externo capaz de
reportar, en un \'unico paso, si una cadena $w$ pertenece o no a $B$. Una
\emph{MT con or\'aculo} $M^{B}$ es una MT modificada con la posibilidad de
hacer consultas a tal or\'aculo. Formalmente se modela con un par de
estados especiales $q_{\text{ask}},q_{\text{yes}},q_{\text{no}}$ y una cinta
adicional para escribir la consulta.
\end{definicion}

\begin{definicion}[Sipser 6.20, Turing reducibility]
$A$ es \emph{Turing reducible} a $B$, notado $A\redT B$, si $A$ es decidible
relativo a $B$; es decir, existe una MT con or\'aculo $M^{B}$ que decide
$A$.
\end{definicion}

\subsection{Propiedades b\'asicas}

\begin{teorema}[Sipser 6.21]
Si $A\redT B$ y $B$ es decidible, entonces $A$ es decidible.
\end{teorema}

\begin{proof}
Sea $M^{B}$ un decisor relativo a $B$ y $D_B$ un decisor (real) de $B$. Una
MT estandar $M'$ simula a $M^{B}$ reemplazando cada consulta al or\'aculo por
una simulacion completa de $D_B$. $M'$ decide $A$ sin or\'aculos.
\end{proof}

\begin{observacion}
\textbf{$\redT$ generaliza a $\redm$.} Si $A\redm B$ con reduccion $f$,
construimos una MT con or\'aculo $B$ que dada $x$ computa $f(x)$ y consulta
``$f(x)\in B$?''. La respuesta del or\'aculo es la respuesta para
``$x\in A$?''. Luego toda mapping reduction es una Turing reduction.
\end{observacion}

\subsection{Ejemplo cl\'asico: $\ETM \redT \ATM$}

\begin{ejemplo}[Sipser 6.19]
$\ETM = \{\enc{M} \mid \lang(M)=\emptyset\}$ es decidible relativo a $\ATM$.

\medskip\noindent\textbf{$T^{\ATM} = $ ``Con entrada $\enc{M}$:}
\begin{enumerate}
  \item construir la MT $N=$ ``En cualquier entrada: ejecutar $M$ en
        paralelo sobre todas las cadenas de $\Sstar$; aceptar si alguna es
        aceptada por $M$'';
  \item consultar al or\'aculo si $\enc{N,0}\in\ATM$;
  \item si la respuesta es NO, aceptar; si es S\'I, rechazar.''
\end{enumerate}
\medskip

\noindent Si $\lang(M)\neq\emptyset$, $N$ acepta toda entrada y la consulta
da S\'I; $T$ rechaza. Si $\lang(M)=\emptyset$, $N$ nunca acepta y la
consulta da NO; $T$ acepta. Por lo tanto $T^{\ATM}$ decide $\ETM$.
\end{ejemplo}

\subsection{L\'imite del poder de los or\'aculos: argumento de cardinalidad}

A\'un teniendo un or\'aculo poderoso, existen lenguajes inalcanzables. La
clave es contar.

\begin{teorema}[Sipser, Problema 6.19]
\label{teo:cardA_TM}
Existen lenguajes $L\subseteq\Sstar$ que no son reconocidos por ninguna MT
con or\'aculo $\ATM$.
\end{teorema}

\begin{proof}
La familia $\mathcal{M}_{\ATM}$ de MT con or\'aculo $\ATM$ es contable
(cada una se describe por una cadena finita). Por lo tanto el conjunto de
lenguajes reconocidos por elementos de $\mathcal{M}_{\ATM}$ es contable. En
cambio, $\Powerset(\Sstar)$ es incontable (Corolario 4.18). Por
comparaci\'on de cardinalidades existe alg\'un $L\in\Powerset(\Sstar)$ fuera
de esa familia.
\end{proof}

\begin{observacion}
El argumento es \emph{robusto bajo cambios de or\'aculo}: para cualquier
$B\subseteq\Sstar$ existen lenguajes no reconocibles por MT con or\'aculo
$B$. Esto da una idea de la \emph{jerarqu\'ia aritm\'etica} y de los grados
de Turing.
\end{observacion}

\bigskip\hrule\bigskip

% =============================================================================
\section{Conexi\'on con el pr\'actico}

\subsection*{Mapa ejercicio $\leftrightarrow$ teor\'ia}

\begin{center}
\begin{tabular}{cll}
\toprule
Ej. & Tema & Herramienta te\'orica principal \\
\midrule
2  & quine en Python                      & Lema~\ref{lem:q} y $\SELF$ \\
3  & dos MT que se imprimen mutuamente    & Lema~\ref{lem:q} + Teorema~\ref{teo:recursion} \\
4  & lenguajes no reconocibles con $M^{\ATM}$ & Teorema~\ref{teo:cardA_TM} \\
5  & introspecci\'on de una clase Python  & \texttt{inspect} como an\'alogo de ``obtener $\enc{M}$'' \\
6  & creaci\'on din\'amica de una clase   & \texttt{type} como an\'alogo de $q$ \\
\bottomrule
\end{tabular}
\end{center}

\subsection*{Resumen ejecutivo de la t\'ecnica de autoreferencia}

\begin{enumerate}
  \item Si necesitamos una MT que use $\enc{M}$, escribimos su pseudoc\'odigo
        usando libremente la frase ``obtener $\enc{M}$'' y aplicamos el
        Teorema~\ref{teo:recursion} para garantizar su existencia.
  \item Si necesitamos una MT que imprima cierto texto fijo $w$, usamos
        $\Pw{w}$ del Lema~\ref{lem:q}.
  \item Si necesitamos una MT que produzca otra MT a partir de su propia
        descripci\'on, combinamos ambos: $\SELF$ + transformaci\'on
        computable $\tau$.
\end{enumerate}

\subsection*{Pseudoc\'odigo unificado}

\begin{verbatim}
M = "En entrada x:
  1. (Si se necesita) obtener <M> por el Teorema de la Recursion.
  2. Computar f(<M>, x) (o cualquier funcion computable usando <M>).
  3. Imprimir, simular, aceptar/rechazar, etc."
\end{verbatim}

\bigskip\hrule\bigskip

% =============================================================================
\section*{S\'intesis final}
\addcontentsline{toc}{section}{S\'intesis final}

El Cap\'itulo 6 de Sipser y las l\'aminas \texttt{autoref.pdf} introducen tres
ideas que se entrelazan:

\begin{itemize}
  \item \textbf{Autoreferencia}: una MT puede obtener su propia descripci\'on
        (Teorema de la Recursi\'on). De aqu\'i se derivan, mec\'anicamente,
        $\SELF$, los quines, la no reconocibilidad de $\MIN$ y una nueva
        prueba de indecibilidad de $\ATM$.
  \item \textbf{Or\'aculos y $\redT$}: extender una MT con un or\'aculo es la
        forma m\'as general de la noci\'on intuitiva de reducci\'on. Sin
        embargo, agregar un or\'aculo no permite escapar de los l\'imites
        cardinales: existen lenguajes incontables y or\'aculos numerables.
  \item \textbf{Introspecci\'on}: en lenguajes Turing-completos reales (Python
        es el ejemplo can\'onico), las facilidades de \texttt{inspect} y
        \texttt{type} hacen tangibles las construcciones de la teor\'ia.
\end{itemize}

En conjunto, este pr\'actico convierte un tema aparentemente filos\'ofico
(``\!\!?puede un programa hablar de s\'i mismo?'') en una herramienta
operativa de la teor\'ia de la computaci\'on y de la programaci\'on real.

\end{document}
